\newpage
\begin{minipage}{0.8\textwidth}
\textbf{Title:} Virtual Reality Prism Effect\\
\textbf{Project:} 8. Semester project\\
\textbf{Project period:} February 2026 -\\ May 2026\\
\end{minipage}
\begin{minipage}{1\textwidth}
\hspace{-1.8cm}
\includegraphics[scale=1]{img/aau_logo_en.pdf}
\end{minipage}\newline\newline

\vspace{5mm}
\begin{minipage}{0.8\textwidth}
\Large\textbf{Project group:} \vspace{5mm}
%\large\textbf{Deltagere:}\\
%\vspace{2mm}

    
    %Person 1%
    \begin{minipage}[b]{0.45\textwidth}
     \centering
     \textit{{\vspace{-6mm} Markus Birch Flensborg}}
     \rule{\textwidth}{0.5pt}\\
     {\footnotesize mflens22@student.aau.dk}
     \vspace{0.5cm}
    \end{minipage}
    
    %Person 2%
    \begin{minipage}[b]{0.45\textwidth}
     \centering
     \textit{{\vspace{-6mm} Mathias Bisgaard}}
     \rule{\textwidth}{0.5pt}\\
     {\footnotesize mbisga22@student.aau.dk}
     \vspace{0.5cm}
    \end{minipage}
    
    %Person 3%
    \begin{minipage}[b]{0.45\textwidth}
     \centering
     \textit{{\vspace{-6mm} Rikke Bragh Jensen}}
     \rule{\textwidth}{0.5pt}\\
     {\footnotesize rbje22@student.aau.dk}
     \vspace{0.5cm}
    \end{minipage}

    %Person 4%
    \begin{minipage}[b]{0.45\textwidth}
     \centering
     \textit{{\vspace{-6mm} Tze Huo Gucci Ho}}
     \rule{\textwidth}{0.5pt}\\
     {\footnotesize tho22@student.aau.dk}
     \vspace{0.5cm}
    \end{minipage}

    
\end{minipage}
\hspace{-5.5cm} \begin{minipage}{0.53\linewidth} \vspace{-1 cm}
\Large\textbf{Abstract:} \vspace{3mm}\\
\noindent\fbox{%
    \parbox{1\textwidth}{% 
\normalsize 

Unilateral spatial neglect is a disorder in which patients fail to attend or respond to one side of space after brain injury. Prism adaptation is a rehabilitation approach where repeated pointing under a visual displacement can produce error reduction during exposure and an aftereffect after the displacement is removed. Virtual reality makes it possible to simulate prism-like perturbations in software, but existing VR approaches differ in what is shifted, how feedback is presented, and which tasks are used to assess adaptation.

This project developed a configurable Unity-based VR sandbox for comparing several implementations of a prism effect within one shared baseline--exposure--post workflow. The prototype implements translation, rotation, and skew perturbations, supports controller and embodied hand-tracking input, records runtime data to CSV, and provides an RShiny dashboard for condition-level, participant-level, quality, spatial, and trajectory analysis.

The prototype was evaluated with 10 healthy participants across 80 completed modules. The evaluation showed that the system could collect analyzable data across tasks, input modes, and perturbation types. However, the results were noisy, and the no-effect condition often produced non-zero baseline-to-post shifts. The clearest condition-level pattern appeared in controller-based Line Bisection, where perturbation conditions showed positive median shifts above participant-specific no-effect references. Participant behaviour, fatigue, tracking stability, posture, confirmation method, exposure quality, and possible carry-over between modules all affected interpretation.

The project therefore contributes primarily as a framework-building and methodological prototype rather than as evidence of clinical rehabilitation efficacy. It shows that VR prism adaptation should be described through explicit dimensions such as perturbation class, exposure visibility, input representation, confirmation method, task family, reset procedure, and analysis method.

}%
}
\end{minipage}\newline\newline\newline\newline

%\vfill
\begin{minipage}{0.5\textwidth}
\textbf{Supervisors:} Bastian Ilsø Hougaard \\
\textbf{Department:} Department of Architecture, Design and Media Technology \\
\textbf{Due date:} May 28th, 2026\\
\end{minipage}
\begin{minipage}{0.5\textwidth}
\begin{flushright}
\textbf{Number of pages: ..} \\
\textbf{Appendix: ..} \\
\end{flushright}
\end{minipage}

\begin{center}
  
  \textbf{Keywords:} Prism Effect; Prism Adaptation; Virtual Reality; C\#; Unity;   \\

  %\vspace{5mm}
    \textsc{The content of this report is freely available, but publication (with source indication) may only be made in agreement with the authors.}

\end{center}
